\chapter{Experimentálne overenie}
\label{ExperimentalneOverenie}

Táto kapitola prezentuje reálne skúšky navrhnutého systému, ktoré overujú splnenie stanovených cieľov zadania v podmienkach reálnej prevádzky.

\section{Popis fyzického testovania a~testovacieho prostredia}
\label{PopisDemonstracie}

Systém bol nasadený a~overený na reálnej fyzickej inštalácii
\textbf{CHRONOS} -- steampunkovej expozičnej stene s~rozmermi
približne 3\,m\,$\times$\,1{,}5\,m. Inštalácia integruje
dekoratívne hodinové mechanizmy, sústavu ozubených kolies,
analógové budíky, dymostroj a~scénické osvetlenie,
ktorých mechanická konštrukcia je popísaná
v~kapitole~\ref{MechanickaKonstrukcia}.

Do testovania boli zapojené nasledujúce uzly a~zariadenia:
\begin{itemize}
    \item \textbf{Raspberry Pi} s~bežiacim centrálnym procesom
        \texttt{MuseumController}, pripojeným HDMI monitorom
        a~audio výstupom cez 3{,}5\,mm jack.
    \item \textbf{ESP32 Motorický uzol} riadiaci dva jednosmerné
        DC motory cez H-mostíky BTS7960 -- motor hodinového
        mechanizmu a~motor ozubených kolies.
    \item \textbf{ESP32 Reléový uzol} (Waveshare ESP32-S3)
        spínajúci osem kanálov: napájanie dymostroja (230\,V),
        produkciu dymu (auto-off 12\,s), svetlo ohňa,
        Edisonovu žiarovku a~štyri kanály podsvietenia
        LED pásov (12\,V).
    \item \textbf{ESP32 Tlačidlový uzol} ako fyzický
        štartovací vstup pre návštevníka.
\end{itemize}

\section{Testované scenáre a výsledky}
\label{TestovaneScenare}

Overenie stability systému prebiehalo v troch izolovaných rovinách: prevádzka s hardvérom, softvérová výdrž jadra a odolnosť voči sieťovým výpadkom.

\subsection{Scenár 1 -- Komplexná lineárna scéna s hardvérom}
\begin{itemize}
    \item \textbf{Cieľ:} Otestovať plynulosť a asynchrónnu presnosť prechodov definovaných v JSON súbore počas riadenia fyzických aktuátorov priamo v~múzeu.
    \item \textbf{Priebeh:} Systém musel prehrať multimediálny audiovizuálny obsah súčasne s presne načasovaným pohybom motorov a dynamickým zhasínaním svetiel.
    \item \textbf{Výsledok:} Všetky akcie boli spustené v~časoch definovaných v~\textit{timeline}. Komunikácia s~mikrokontrolérmi prebehla bez zaznamenaného oneskorenia a~synchronizácia audio-video stopy so svetelnými efektmi bola vizuálne konzistentná počas celého trvania testu.
\end{itemize}

\subsection{Scenár 2 -- Dlhodobý test stability riadiaceho systému}
\begin{itemize}
    \item \textbf{Cieľ:} Overiť odolnosť Raspberry Pi a MQTT brokera proti pretečeniu pamäte a zlyhaniu pri dlhodobom plynulom behu (tzv. \textit{Soak test}).
    \item \textbf{Priebeh:} Centrálny program systému bežal na stole formou \textit{dry-run} testovania v kuse niekoľko dní v laboratórnych podmienkach. Priebežne dochádzalo k ojedinelému spúšťaniu scén a monitorovaniu odosielaných MQTT príkazov na zbernicu.
    \item \textbf{Výsledok:} Proces \texttt{MuseumController} zostal počas celej doby testu bez reštartu. Príkazy boli generované na MQTT zbernicu priebežne bez výpadku, čo preukázalo dostatočnú stabilitu softvéru pre nepretržitú prevádzku.
\end{itemize}

\subsection{Scenár 3 -- Test sieťovej odolnosti a LWT (Last Will)}
\begin{itemize}
    \item \textbf{Cieľ:} Overiť detekčný a oznamovací mechanizmus systému pri neštandardnej strate bezdrôtového spojenia uzlov na Wi-Fi sieti.
    \item \textbf{Priebeh:} Koncovým ESP32 uzlom (\textit{Controller\_MOTORS} a \textit{Controller\_RELAY}) bolo striedavo fyzicky odpájané a následne vracané napájanie za behu celého systému.
    \item \textbf{Výsledok:} Po obnovení napájania sa uzly autonómne pripojili späť na Wi-Fi sieť a~MQTT brokera. Dashboard správne zaznamenal zmenu stavu - prepis indikátora z~\texttt{offline} na \texttt{online} prebehol na základe LWT správy od brokera a~následného \texttt{online} hlásenia od uzla.
\end{itemize}

% ---------------------------------------------------------------
% Vložiť za \subsection{Scenár 3 -- Test sieťovej odolnosti...}
% a PRED \section{Zhodnotenie overenia}
% ---------------------------------------------------------------
 
\subsection{Scenár 4 -- Meranie latencie MQTT spätnej väzby}
\label{LatenciaMQTT}
\begin{itemize}
    \item \textbf{Cieľ:} Kvantitatívne overiť latenciu cyklu
        príkaz\,--\,spätná väzba medzi backendom Raspberry~Pi
        a~motorickým uzlom ESP32 pri vysokofrekvenčnom zaťažení
        (10 príkazov za sekundu, stres-test TC-08).
    \item \textbf{Priebeh:} Počas behu stres-testovej scény bolo
        z~databázového záznamu modulu \texttt{mqtt\_feedback}
        extrahovaných 128 hodnôt doby odozvy -- po 64 pre každý
        motorický kanál. Meranie pokrývalo súvislý časový úsek
        40 sekúnd.
    \item \textbf{Výsledok:} Priemerná latencia celého cyklu
        dosiahla 88{,}7\,ms pri mediáne 90\,ms. Žiadna
        zo 128 meraných hodnôt neprekročila hranicu 200\,ms --
        59{,}4\,\% meraní spadalo do pásma pod 100\,ms
        a~zvyšných 40{,}6\,\% do pásma 100\,--\,200\,ms
        (Tab.~\ref{tab:mqtt_latency}, Obr.~\ref{fig:mqtt_latency}).
        Hodnota P99 dosiahla 166\,ms, čo potvrdzuje absenciu
        extrémnych odchýlok. Kanál \texttt{motor2} vykazoval
        v~priemere o~17{,}9\,ms vyššiu latenciu ako
        \texttt{motor1}; rozdiel je zhodný s~typickým
        variabilným oneskorením domácej Wi-Fi siete
        a~nemá vplyv na pozorovateľnosť efektov.
\end{itemize}
 
\begin{table}[H]
\centering
\caption{Štatistika latencie MQTT cyklu príkaz\,--\,spätná väzba}
\label{tab:mqtt_latency}
\begin{tabular}{|l|c|c|c|c|c|c|}
\hline
\textbf{Zariadenie} & \textbf{n} & \textbf{Priemer (ms)}
    & \textbf{Medián (ms)} & \textbf{Min (ms)}
    & \textbf{Max (ms)} & \textbf{P95 (ms)} \\
\hline
motor1 & 64 & 79{,}7 & 81{,}0 & 19 & 143 & 136 \\
motor2 & 64 & 97{,}6 & 102{,}0 & 17 & 167 & 164 \\
\hline
\textbf{Celkovo} & \textbf{128} & \textbf{88{,}7}
    & \textbf{90{,}0} & \textbf{17} & \textbf{167}
    & \textbf{156} \\
\hline
\end{tabular}
\end{table}
 
\begin{figure}[H]
\centering
\includegraphics[width=0.9\linewidth]{obr/mqtt_latency}
\caption{Distribúcia latencie MQTT spätnej väzby (n = 128,
    stres-test TC-08, 10\,Hz)}
\label{fig:mqtt_latency}
\end{figure}

\section{Zhodnotenie overenia}
\label{VysledkyZhodnotenie}

Všetky tri testované scenáre boli úspešne overené. Architektúra
distribuovanej hviezdicovej siete sa ukázala ako dostatočne stabilná
pre podmienky plného fyzického zaťaženia v~prostredí expozície.
Bezdrôtová komunikácia cez protokol MQTT preukázala reaktivitu
postačujúcu na spracovanie hardvérových triggerov bez merateľného
oneskorenia, čo vylúčilo potrebu prechodu na drôtové prepojenie.
